\chapterHeading{Miscellaneous}


\subHeading{Statement notes}
\begin{itemize}
\item leap year: \texttt{(!(y\%4) \&\& (y\%100)) || !(y\%400)}
\item Harmonic series generates at most $2\sqrt{N}$ distinct values of $\lfloor N/i \rfloor$.
\item Prime gap: $\log^2 N$.
\item Sometimes from any index $i$, the cost is (cost from 1 to destination $-$ destination distance).
\item \strong{mean} works on Euclidean distance; \strong{median} works on Manhattan distance.
\item To avoid TLE with \texttt{map}/\texttt{unordered\_map}:
\begin{lstlisting}[language=C++]
unordered_map<int, int> mp;
mp.reserve(n * n);
\end{lstlisting}
\item Combine a pair into a single hash value:
\begin{lstlisting}[language=C++]
// both non-negative:
   ((uint64_t)p.first << 32) | p.second;
// if a negative may occur:
   ((int)x << 32) | (y & 0xffffffffLL);
\end{lstlisting}
\end{itemize}

\subHeading{Adhoc}
\begin{itemize}
\item prefix value / suffix value / prefix $+$ suffix value
\item print all answers for 1 to 100 and look for a pattern
\item check the equation by value
\item check units for equation validity
\item symmetric cases
\item the values don't matter, the order does
\end{itemize}

\subHeading{Bit tricks}
\begin{lstlisting}[language=C++]
__builtin_popcountll(n);  // number of set bits
__builtin_clzll(x);       // leading zeros
__builtin_ctz(x);         // trailing zeros
__builtin_ffs(x);         // index of first set bit
__lg(x);                  // highest set bit (n != 0)
bit.to_ullong();          // bitset -> unsigned long long

((n & (1 << i)) != 0);    // i-th bit set?
n = (n | (1 << i));       // set i-th bit
n = (n & (~(1 << i)));    // clear i-th bit
n = (n ^ (1 << i));       // toggle i-th bit
n = (n & (n - 1));        // clear rightmost set bit
x = !5;                   // x = -6

// n         = 00110100
// n-1       = 00110011
// n & (n-1) = 00110000
\end{lstlisting}
The number of set bits among numbers up to $2^x$ is $x \cdot 2^{x-1}$.

\subHeading{Useful functions}
\begin{lstlisting}[language=C++]
expl(x); // e^x		also logl, log10l, log2l, cbrtl
// remove duplicates (sorted)
   vec.erase(unique(all(vec)), vec.end());   
double (*\codebf{log_x_of_base_b}*)(int x, int b)
   return log(x)/log(b);	// log of x in base b
// convert x from base b1 to base b2
double (*\codebf{base_b1_to_b2}*)(int x, int b1, int b2)
     return log_x_of_base_b(x, b1)/log_x_of_base_b(b2, b1);
// digit sum until a single digit
cout << (n % 9 ? n % 9 : n ? 9 : 0) << nl;
//parse in any base
int x = (*\codebf{stoll}*)(binaryNumber, nullptr, 2); 
// leap years in (l, r]
int (*\codebf{LeapYearInRange}*)(int r, int l) {        
    l--;
    int R = r / 4 - r / 100 + r / 400;
    int L = l / 4 - l / 100 + l / 400;
    return R - L;
}
\end{lstlisting}

\subHeading{Principle of Inclusion and Exclusion}
\[|A \cup B \cup C| = |A| + |B| + |C| - |A \cap B| - |A \cap C| - |B \cap C| + |A \cap B \cap C|\]
\begin{lstlisting}[language=C++]
int PIE(vector<int>& v, int n, int m) {   // O(2^n * n)
    int even = 0, odd = 0;
    for (int mask = 1; mask < (1LL << n); mask++) {
        int tmp = m;
        for (int i = 0; i < n; i++)
            if (mask & (1LL << i)) tmp /= v[i];
        if (__builtin_popcountll(mask) & 1) odd += tmp;
        else even += tmp;
    }
    return odd - even;
}
\end{lstlisting}

\subHeading{LCS via LIS}
\begin{lstlisting}[language=C++]
// LCS -> LIS:
// Store the positions of each value of A in B.
// For each A[i], add its B positions to cs in REVERSE order
//   (this handles duplicates).
// LCS length = LIS(cs).
// If each value occurs <= K times: |cs| <= N*K
// TC: O(N*K log(N*K)); worst O(N^2 log N) if one value occurs n times.
\end{lstlisting}

\subHeading{Ternary search}
\begin{lstlisting}[language=C++]
m1 = lo + (hi - lo) / 3.0;
m2 = hi - (hi - lo) / 3.0;
\end{lstlisting}

\iffalse
\sectionHeading{Matrix}
\begin{lstlisting}[language=C++]
struct Matrix {
    vector<vector<int>> v;
    int get(int r, int c) { return v[r][c]; }
    int rowCount() { return v.size(); }
    int colCount() { return v[0].size(); }
};

Matrix operator*(Matrix a, Matrix b) {
    Matrix result;
    int p = a.rowCount(), q = a.colCount();   // q = b.rowCount()
    int r = b.colCount();
    result.v.resize(p);
    for (int i = 0; i < p; i++) result.v[i].resize(r);
    for (int i = 0; i < p; i++)
        for (int j = 0; j < q; j++)
            for (int k = 0; k < r; k++)
                result.v[i][k] += a.get(i, j) * b.get(j, k);
    return result;
}

Matrix MatExpo(Matrix X, int p) {
    Matrix result;
    int r = X.rowCount();
    result.v.resize(r);
    for (int i = 0; i < r; i++) result.v[i].resize(r);
    for (int i = 0; i < r; i++) result.v[i][i] = 1;
    while (p) {
        if (p & 1) result = result * X;
        X = X * X;
        p >>= 1;
    }
    return result;
}
\end{lstlisting}
\fi
